Consideraremos que el tope de la pila está a la derecha, por lo que sacamos y evaluamos este último elemento.

\begin{enumerate}[label=\arabic*:]
\item Metemos el estado inicial \texttt{(1,2)} a la pila, además de marcarlo como visitado:
  \begin{align*}
    \text{Pila:} &\quad \texttt{[(1,2)]} \\
    \text{Visitados:} &\quad \{(1,2)\}
  \end{align*}

  Se genera la raíz del árbol:
  \begin{figure}[h]
    \centering
    \begin{forest}
      for tree={ math content, edge={-}, s sep=25pt, l sep=20pt, align=center }
      [{(1,2)}]
    \end{forest}
  \end{figure}
  
\item Sacamos \texttt{(1,2)} y evaluamos sus vecinos:
  
  \textit{Izquierda} \texttt{(1,1)}, \textit{Derecha} \texttt{(1,3)} (\textit{Arriba} y \textit{Abajo} no son válidos). Metemos estas acciones en orden inverso para así aplicarlas en el orden que se pide, de modo que la \textit{Izquierda} se mete al final, lo que la hace estar en el tope de la pila y sea la primer acción ejecutada:
  \begin{align*}
    \text{Pila:} &\quad \texttt{[(1,3), (1,1)]} \\
    \text{Visitados:} &\quad \{(1,2), (1,1), (1,3)\}
  \end{align*}

  El árbol se expande:
  \begin{figure}[h]
    \centering
    \begin{forest}
      for tree={ math content, edge={-}, s sep=25pt, l sep=20pt, align=center }
      [{(1,2)}
        [{(1,1)}, edge label={node[midway,left]{$\leftarrow$}}]
        [{(1,3)}, edge label={node[midway,right]{$\rightarrow$}}]
      ]
    \end{forest}
  \end{figure}
  
\item Sacamos \texttt{(1,1)}. Metemos \textit{Abajo} \texttt{(2,1)}:
  \begin{align*}
    \text{Pila:} &\quad \texttt{[(1,3), (2,1)]} \\
    \text{Visitados:} &\quad \{(1,2), (1,1), (1,3), (2,1)\}
  \end{align*}

  De este modo nos vamos extendiendo en el árbol:
  \begin{figure}[h]
    \centering
    \begin{forest}
      for tree={ math content, edge={-}, s sep=25pt, l sep=20pt, align=center }
      [{(1,2)}
        [{(1,1)}, edge label={node[midway,left]{$\leftarrow$}}]
        [{(1,3)}, edge label={node[midway,right]{$\rightarrow$}}]
      ]
    \end{forest}
  \end{figure}

  \newpage
  
\item Sacamos \texttt{(2,1)}. Metemos \textit{Abajo} \texttt{(3,1)}:
  \begin{align*}
    \text{Pila:} &\quad \texttt{[(1,3), (3,1)]} \\
    \text{Visitados:} &\quad \{(1,2), (1,1), (1,3), (2,1), (3,1)\}
  \end{align*}

  Nos seguimos extendiendo por el nodo:
  \begin{figure}[h]
    \centering
    \begin{forest}
      for tree={ math content, edge={-}, s sep=25pt, l sep=20pt, align=center }
      [{(1,2)}
        [{(1,1)}, edge label={node[midway,left]{$\leftarrow$}}
          [{(2,1)}, edge label={node[midway,left]{$\downarrow$}}]
        ]
        [{(1,3)}, edge label={node[midway,right]{$\rightarrow$}}]
      ]
    \end{forest}
  \end{figure}
  
\item Sacamos \texttt{(3,1)}. No hay acciones válidas, por lo que solo queda la pila:
  \begin{align*}
    \text{Pila:} &\quad \texttt{[(1,3)]}
  \end{align*}

  Y se termina la extensión del árbol por esta rama:

  \begin{figure}[h]
    \centering
    \begin{forest}
      for tree={ math content, edge={-}, s sep=25pt, l sep=20pt, align=center }
      [{(1,2)}
        [{(1,1)}, edge label={node[midway,left]{$\leftarrow$}}
          [{(2,1)}, edge label={node[midway,left]{$\downarrow$}}
            [{(3,1)}, edge label={node[midway,left]{$\downarrow$}}, label={[yshift=-2pt]below:{\Large \color{red}$\times$}}]
          ]
        ]
        [{(1,3)}, edge label={node[midway,right]{$\rightarrow$}}]
      ]
    \end{forest}
  \end{figure}
  
  Continuamos extendiéndonos por el otro nodo \texttt{(1,3)}.
  
\item Sacamos \texttt{(1,3)}. Metemos a sus vecinos \textit{Derecha} \texttt{(1,4)}, \textit{Abajo} \texttt{(2,3)}; en ese orden (para evaluar \textit{Abajo} primero):
  \begin{align*}
    \text{Pila:} &\quad \texttt{[(1,4), (2,3)]} \\
    \text{Visitados:} &\quad \{(1,2), (1,1), (1,3), (2,1), (3,1), (2,3), (1,4)\}
  \end{align*}
  
  Y se forma el árbol:

  \begin{figure}[h]
    \centering
    \begin{forest}
      for tree={ math content, edge={-}, s sep=25pt, l sep=20pt, align=center }
      [{(1,2)}
        [{(1,1)}, edge label={node[midway,left]{$\leftarrow$}}
          [{(2,1)}, edge label={node[midway,left]{$\downarrow$}}
            [{(3,1)}, edge label={node[midway,left]{$\downarrow$}}, label={[yshift=-2pt]below:{\Large \color{red}$\times$}}]
          ]
        ]
        [{(1,3)}, edge label={node[midway,right]{$\rightarrow$}}
          [{(2,3)}, edge label={node[midway,left]{$\downarrow$}}]
          [{(1,4)}, edge label={node[midway,right]{$\rightarrow$}}]
        ]
      ]
    \end{forest}
  \end{figure}
  
  Nos extenderemos en el nodo \texttt{(2,3)}.
  
\item Sacamos \texttt{(2,3)}. Metemos \textit{Abajo} \texttt{(3,3)}:
  \begin{align*}
    \text{Pila:} &\quad \texttt{[(1,4), (3,3)]} \\
    \text{Visitados:} &\quad \{(1,2), (1,1), (1,3), (2,1), (3,1), (2,3), (1,4), (3,3)\}
  \end{align*}
  
\item Sacamos \texttt{(3,3)}. Metemos \textit{Derecha} \texttt{(3,4)}:
  \begin{align*}
    \text{Pila:} &\quad \texttt{[(1,4), (3,4)]} \\
    \text{Visitados:} &\quad \{(1,2), (1,1), (1,3), (2,1), (3,1),(2,3), (1,4), (3,3), (3,4)\}
  \end{align*}
  
\item Sacamos \texttt{(3,4)}. Metemos \textit{Abajo} \texttt{(4,4)} que es la meta, y por usar \textbf{Early Goal Test} terminamos el algoritmo inmediatamente, pues ya encontramos la solución.
\end{enumerate}

\newpage

\textbf{Árbol de búsqueda completo:}
\begin{figure}[h]
  \centering
  \begin{forest}
    for tree={ math content, edge={-}, s sep=25pt, l sep=20pt, align=center }
    [{(1,2)}
      [{(1,1)}, edge label={node[midway,left]{$\leftarrow$}}
        [{(2,1)}, edge label={node[midway,left]{$\downarrow$}}
          [{(3,1)}, edge label={node[midway,left]{$\downarrow$}}, label={[yshift=-2pt]below:{\Large \color{red}$\times$}}]
        ]
      ]
      [{(1,3)}, edge label={node[midway,right]{$\rightarrow$}}
        [{(2,3)}, edge label={node[midway,left]{$\downarrow$}}
          [{(3,3)}, edge label={node[midway,left]{$\downarrow$}}
            [{(3,4)}, edge label={node[midway,left]{$\rightarrow$}}
              [{(4,4)}, edge label={node[midway,left]{$\downarrow$}}, label={[yshift=-2pt]below:{\Large \color{green}$\checkmark$}}]
            ]
          ]
        ]
        [{(1,4)}, edge label={node[midway,right]{$\rightarrow$}}]
      ]
    ]
  \end{forest}
\end{figure}

Cabe aclarar que al usar DFS, nunca exploramos el nodo \texttt{(1,4)}. Si no invertimos la forma de agregar las acciones a la pila el árbol en cambio sería:

\newpage

\begin{figure}[h]
  \centering
  \begin{forest}
    for tree={ math content, edge={-}, s sep=25pt, l sep=20pt, align=center }
    [{(1,2)}
      [{(1,1)}, edge label={node[midway,left]{$\leftarrow$}}]
      [{(1,3)}, edge label={node[midway,right]{$\rightarrow$}}
        [{(2,3)}, edge label={node[midway,left]{$\downarrow$}}
          [{(3,3)}, edge label={node[midway,left]{$\downarrow$}}
            [{(3,4)}, edge label={node[midway,left]{$\rightarrow$}}
              [{(4,4)}, edge label={node[midway,left]{$\downarrow$}}, label={[yshift=-2pt]below:{\Large \color{green}$\checkmark$}}]
            ]
          ]
        ]
        [{(1,4)}, edge label={node[midway,right]{$\rightarrow$}}]
      ]
    ]
  \end{forest}
\end{figure}
